%%
%% This is file `min-exemplo.tex',
%% generated with the docstrip utility.
%%
%% The original source files were:
%%
%% ufrj-coppe.dtx  (with options: `minexemplo')
%% 
%% This is a sample monograph which illustrates the use of the `ufrj' document
%% class with the `ufrj-coppe' unit style and its biblatex/biber bibliography
%% style. It is generated from ufrj-coppe.dtx by docstrip.
%% 
%% Copyright (C) 2008-2026 Vicente Helano F. Batista, George O. Ainsworth Jr.,
%% Geraldo Xexéo and Eduardo Mangeli. Released under the GNU General Public
%% License version 3 (see the COPYING.txt file).
%% 
%% !TeX program = pdflatex
%% !TeX encoding = UTF-8
%% !TeX spellcheck = pt_BR
%% !BIB program = biber
%% Com o latexmk funcionando (no Windows ele pede o Perl), troque a primeira
%% linha por: % !TeX program = latexmk -pdf -synctex=1
\documentclass[dsc]{ufrj}
\usepackage{ufrj-coppe}
%% min-exemplo.tex -- tudo o que a norma exige, e nada além disso.
%%
%% O que está aqui e por quê (Manual UFRJ/SiBI, 9.ª ed. rev., 2026):
%%  - PDF/A ............... o depósito é em PDF/A (2.2d), e é o padrão da
%%    classe: não se escreve opção nenhuma para ter;
%%  - \maketitle .......... capa, folha de rosto (3.1.2.1.1) e folha adicional
%%    com a Coleta CAPES e a ficha catalográfica (3.1.2.1.2);
%%  - \frontmatter ........ folha de aprovação (3.1.2.1.3), com o orientador
%%    em primeiro lugar na banca, como presidente;
%%  - abstract ............ resumo na língua do trabalho (3.1.2.1.4);
%%  - foreignabstract ..... resumo em língua estrangeira (3.1.2.1.5);
%%  - \tableofcontents .... sumário (3.1.2.1.6), o único obrigatório das listas;
%%  - introdução, desenvolvimento e conclusão (3.1.3);
%%  - \printbibliography .. referências, o único pós-textual obrigatório (3.1.4).
%%
%% O que NÃO está, por ser opcional: dedicatória, agradecimentos, epígrafe,
%% todas as listas além do sumário, coorientador, linha de pesquisa,
%% glossário, apêndices, anexos e índice. Tudo isso está no max-exemplo.tex.

\addbibresource{exemplo.bib}% a base de referências

\begin{document}
  \title{Título do Trabalho}
  \foreigntitle{Title of the Work}
  % O subtítulo é opcional (3.1.2.1.1c); está aqui só para dar nome ao exemplo.
  % Sai depois do título, com dois-pontos, e sem as maiúsculas do título.
  \subtitle{um exemplo mínimo de uso do CoppeTeX}
  \foreignsubtitle{a minimal example of use of CoppeTeX}
  \author{Nome do Autor}{Sobrenome}
  \advisor{Nome do Orientador}{Sobrenome}{D.Sc.}{UFRJ}
  % A composição da banca é regra do Programa, não da formatação.
  \examiner{Nome do Primeiro Examinador Sobrenome}{D.Sc.}{UFRJ}
  \examiner{Nome do Segundo Examinador Sobrenome}{D.Sc.}{UFRJ}
  \examiner{Nome do Terceiro Examinador Sobrenome}{D.Sc.}{UFF}
  \department{PESC}
  \date{09}{2026}
  % Data da aprovação: obrigatória na folha de aprovação (3.1.2.1.3d).
  \approvaldate{15 de setembro de 2026}

  % Palavras-chave, ao final de cada resumo (3.1.2.1.4 e 3.1.2.1.5).
  \keyword{primeira palavra-chave}
  \keyword{segunda palavra-chave}
  \keyword{terceira palavra-chave}
  \foreignkeyword{First keyword}
  \foreignkeyword{Second keyword}
  \foreignkeyword{Third keyword}

  % Folha adicional: a área de concentração e os demais campos da Coleta CAPES
  % (3.1.2.1.2). A área não sai na folha de rosto -- é o modelo do SiBI, e a
  % COPPE não liga a opção areanafolhaderosto. A ficha catalográfica vem do
  % gerador do SiBI; com o arquivo em mãos, descomente a linha \catalogcard.
  \concentrationarea{Engenharia de Sistemas e Computação}
  \productiontype{bibliographic}
  \linkedproject{no}
  % \catalogcard{ficha.pdf}

  % A ordem dos elementos pré-textuais é a da 3.1.2 do Manual, e a classe a
  % confere: um elemento fora dela é erro de compilação, e a mensagem diz o
  % que mover. Os obrigatórios são estes, nesta ordem: \maketitle (capa,
  % folha de rosto e folha adicional), \frontmatter (folha de aprovação), o
  % resumo em português, o resumo em língua estrangeira e \tableofcontents
  % (sumário), sempre o último. O max-exemplo.tex mostra onde entram os
  % opcionais: errata, dedicatória, agradecimentos, epígrafe e as listas.
  \maketitle

  \frontmatter

  \begin{abstract}
  Resumo do trabalho, de 150 a 500 palavras, num só parágrafo (3.1.2.1.4).
  Lorem ipsum dolor sit amet, consectetur adipiscing elit, sed do eiusmod
  tempor incididunt ut labore et dolore magna aliqua.
  \end{abstract}

  \begin{foreignabstract}
  Abstract of the work, in English (3.1.2.1.5). Lorem ipsum dolor sit amet,
  consectetur adipiscing elit, sed do eiusmod tempor incididunt ut labore et
  dolore magna aliqua.
  \end{foreignabstract}

  \tableofcontents

  \mainmatter

  \chapter{Introdução}

  A introdução apresenta o tema, o problema e os objetivos do trabalho, e
  anuncia como o texto está organizado. Toda afirmação que venha de outro
  trabalho é citada, como em \citet{book-example}.

  \chapter{Desenvolvimento}

  O desenvolvimento pode ocupar quantos capítulos o trabalho pedir: revisão
  da literatura, métodos, resultados e discussão \citep{article-example}.

  \chapter{Conclusão}

  A conclusão responde aos objetivos postos na introdução.

  \backmatter

  \printbibliography

\end{document}
%% 
%%
%% End of file `min-exemplo.tex'.
